\documentclass[12pt, a4paper]{article}


\usepackage[utf8]{inputenc}
\usepackage[italian]{babel}
\usepackage[T1]{fontenc}
\usepackage{lmodern}
\usepackage{amsmath, amssymb}
\usepackage{geometry}
\usepackage{enumitem}
\usepackage{booktabs}
\usepackage{eurosym}
\usepackage{graphicx}
\usepackage{titlesec}
\usepackage{tcolorbox}
\usepackage{hyphenat}
\usepackage{tocloft}
\usepackage{xcolor}
\usepackage{hyperref}

\usepackage{datetime} 
\usepackage{array}
\usepackage{float}

% --- AGGIUNTI PER IL D1 DI TEPO ---
\usepackage{listings}     % per i blocchi PlantUML in UD.tex e UF.tex
\usepackage{longtable}    % per le tabelle di Traceability.tex e GestioneModifiche.tex

% --- CONFIGURAZIONE LISTINGS ---
\definecolor{codegray}{rgb}{0.5,0.5,0.5}
\definecolor{codegreen}{rgb}{0.1,0.5,0.1}
\definecolor{codeblue}{rgb}{0.1,0.1,0.6}
\lstset{
  basicstyle=\ttfamily\small,
  keywordstyle=\color{codeblue}\bfseries,
  commentstyle=\color{codegreen}\itshape,
  stringstyle=\color{codegray},
  frame=single,
  breaklines=true,
  showstringspaces=false,
  captionpos=b,
  columns=fullflexible,
  keepspaces=true
}

% Define \RF for functional requirements

\makeatletter
% RF
\def\toclevel@RF{2}
\def\toclevel@subRF{3}

% RNF
\def\toclevel@RNF{2}
\def\toclevel@subRNF{3}

% User Stories
\def\toclevel@US{2}
\def\toclevel@UF{2}

% --- Section format ---
\titleformat{\section}
  {\normalfont\LARGE\bfseries}  % Font style: normal + Large + bold
  {\thesection}                 % Section number format
  {6pt}                         % Space between number and title
  {}

% ---------------------------
% REQUISITI FUNZIONALI (RF)
% ---------------------------
\newcounter{reqfunc}                  % contatore RF
\newcounter{subRFcounter}[reqfunc]    % contatore sotto-RF

% Dichiaro due livelli personalizzati per RF e subRF
\makeatletter
\newcommand{\l@RF}{\@dottedtocline{2}{1.5em}{2.3em}}
\newcommand{\l@subRF}{\@dottedtocline{3}{3.8em}{3.4em}}
\makeatother

\renewcommand{\thesubRFcounter}{\arabic{reqfunc}.\arabic{subRFcounter}}

\newcommand{\RF}[1]{%
  \refstepcounter{reqfunc}%
  \setcounter{subRFcounter}{0}%

  \addcontentsline{toc}{RF}{RF \arabic{reqfunc}: #1}%
  
  \paragraph*{RF \arabic{reqfunc}: #1}%
}

\newcommand{\subRF}[1]{%
  \refstepcounter{subRFcounter}%
  \addcontentsline{toc}{subRF}{RF \thesubRFcounter: #1}%
  \noindent\hspace{1.5em}\textbf{RF \thesubRFcounter:} #1\par
}

% ---------------------------
% REQUISITI NON FUNZIONALI (RNF)
% ---------------------------

\newcounter{reqnonfunc}
\newcounter{subRNFcounter}[reqnonfunc]

% Livelli personalizzati per TOC (identici a RF)
\makeatletter
\newcommand{\l@RNF}{\@dottedtocline{2}{1.5em}{2.3em}}
\newcommand{\l@subRNF}{\@dottedtocline{3}{3.8em}{3.4em}}
\makeatother

\renewcommand{\thesubRNFcounter}{\arabic{reqnonfunc}.\arabic{subRNFcounter}}

\newcommand{\RNF}[1]{%
  \refstepcounter{reqnonfunc}%
  \setcounter{subRNFcounter}{0}%
  \addcontentsline{toc}{RNF}{RNF \arabic{reqnonfunc}: #1}%
  \paragraph*{RNF \arabic{reqnonfunc}: #1}%
}

\newcommand{\subRNF}[1]{%
  \refstepcounter{subRNFcounter}%
  \addcontentsline{toc}{subRNF}{RNF \thesubRNFcounter: #1}%
  \noindent\hspace{1.5em}\textbf{RNF \thesubRNFcounter:} #1\par
}

% ---------------------------
% User Stories (US)
% ---------------------------

\newcounter{userStories}

% Livelli personalizzati per TOC (identici a RF)
\makeatletter
\newcommand{\l@US}{\@dottedtocline{2}{1.5em}{2.3em}}

\makeatother

\newcommand{\US}[1]{%
  \refstepcounter{userStories}%
  \addcontentsline{toc}{US}{US \arabic{userStories}: #1}%
  \subsection*{User story \arabic{userStories}: #1}%
}

% ---------------------------
% User Flow (UF)
% ---------------------------

\newcounter{userFlows}

% Livelli personalizzati per TOC (identici a RF)
\makeatletter
\newcommand{\l@UF}{\@dottedtocline{2}{1.5em}{2.3em}}

\makeatother

\newcommand{\UF}[1]{%
  \refstepcounter{userFlows}%
  \addcontentsline{toc}{UF}{UF \arabic{userFlows}: #1}%
  \subsection*{User flow \arabic{userFlows}: #1}%
} 

\newcommand{\myfootertext}{}
\newcommand{\myicon}{logo_unitn_chiaro.png}
\newcommand{\myappicon}{logo_app.png}

\definecolor{lightergray}{gray}{0.35}
\definecolor{tocsec}{gray}{0.35}
\definecolor{tocsubsec}{gray}{0.45}
\pagecolor{white}
\color{black}


\hypersetup{
    colorlinks=true,
    linkcolor=lightergray,
    urlcolor=cyan,
    pdftitle={TEPO - D1},
    pdfauthor={Gruppo T14}
}

\renewcommand{\cftsecfont}{\color{tocsec}\bfseries}
\renewcommand{\cftsecpagefont}{\color{tocsec}\bfseries}
\renewcommand{\cftsubsecfont}{\color{tocsubsec}}
\renewcommand{\cftsubsecpagefont}{\color{tocsubsec}}
\renewcommand{\cftsubsecleader}{\color{tocsubsec}\cftdotfill{\cftdotsep}}

\geometry{margin=2.5cm}

\begin{document}

\begin{titlepage}
    % --- Loghi in alto (Sinistra e Destra) ---
    \begin{minipage}[t]{0.45\textwidth}
        \hspace{0cm}
        \href{https://www.unitn.it/}{
            \includegraphics[height=4.2cm, keepaspectratio]{\myicon}%
        }

        \vspace{0pt}
        \href{https://github.com/mirconegri/TEPO}{
            \includegraphics[height=4.2cm, keepaspectratio]{\myappicon}
        }
    \end{minipage}

    \vspace{1.5cm}

    % --- Titolo e Dipartimento (Allineati a sinistra) ---
    \begin{flushleft}
        {\scshape\Large Università degli Studi di Trento \par}
        {\large Dipartimento di Ingegneria e Scienza dell'Informazione \par}
        
        \vspace{1.5cm}
        
        {\Huge\bfseries TEPO - D1 \par}
        \vspace{0.3cm}
        {\Large\itshape Trento Energy Portfolio Optimizer \par}
        \vspace{0.5cm}
        {\color{cyan}\rule{\textwidth}{2pt}}
    \end{flushleft}

    \vspace{4cm}
    
    % --- Gruppo e Autori ---
    \begin{flushleft}
        {\Large \textbf{Gruppo T14} \par}
        \vspace{0.3cm}
        \renewcommand{\arraystretch}{1.3}
        \begin{tabular}{@{}l}
            \Large \href{https://github.com/mirconegri}{Mirco Negri} \\
            \Large \href{https://github.com/ghebbb}{Gabriele Duca} \\
            \Large \href{https://github.com/scheddiee}{Riccardo Scheda} \\
        \end{tabular}
    \end{flushleft}
    
    \vfill
    
    % --- Piè di pagina ---
    \begin{flushleft}
        {\small \textit{\myfootertext} \par}
        {\large \monthname\ \the\year \par}
    \end{flushleft}
\end{titlepage}

\newpage
\begingroup
\hypersetup{hidelinks}
\tableofcontents
\endgroup

\newpage

\include{sections/Il_nostro_progetto}
\section{Requisiti Funzionali}

\subsection{Requisiti funzionali di catalogo e consultazione}

\RF{Consultazione del catalogo degli edifici.} \label{rf:catalogo_edifici}
Il sistema deve consentire agli utenti di consultare il catalogo degli edifici pubblici di interesse, con le informazioni anagrafiche e geometriche essenziali per la valutazione degli interventi (identificativo, denominazione, indirizzo, coordinate geografiche, tipologia d'uso, superficie, anno di costruzione quando disponibile). Per ciascun edificio il sistema deve rendere visibili gli interventi candidati a esso associati. Il sistema deve inoltre rendere esplicita la provenienza del dato (reale, stimato, sintetico) per ciascun edificio.

\RF{Visualizzazione geografica degli edifici.} \label{rf:mappa}
Il sistema deve offrire una vista cartografica degli edifici, con la possibilità di individuare rapidamente la collocazione di ciascun immobile sul territorio comunale. La mappa deve permettere di distinguere visivamente gli edifici che ospitano interventi selezionati, esclusi o non valutati in uno scenario corrente.

\RF{Consultazione del catalogo degli interventi.} \label{rf:catalogo_interventi}
Il sistema deve consentire di consultare l'insieme delle opzioni di intervento disponibili, mostrando per ciascuna: edificio di riferimento, tipologia (LED, serramenti, isolamento, pompa di calore, fotovoltaico, caldaia), costo stimato, risparmio energetico annuo atteso, risparmio economico annuo atteso, riduzione annua attesa di CO\textsubscript{2}, tempo di ritorno semplice e stato di valutazione rispetto a uno scenario corrente.

\RF{Ricerca e filtraggio degli interventi.} \label{rf:filtri_interventi}
Il sistema deve permettere di ricerca e filtrare le opzioni di intervento in base a molteplici criteri combinabili (ricerca testuale, edificio di riferimento, tipologia di intervento, intervallo di costo, stato di valutazione in uno scenario). L'utente deve poter ordinare i risultati per le principali grandezze numeriche (costo, risparmio energetico, risparmio economico, CO\textsubscript{2}, payback).

\RF{Importazione di dati di catalogo.} \label{rf:import_csv}
Il sistema deve consentire l'importazione di dati di catalogo (edifici e opzioni di intervento) da file in formato CSV. L'importazione deve avvenire in due fasi, con una fase di validazione \emph{a secco} che produce un'anteprima delle righe valide, di quelle scartate e delle relative motivazioni, e una fase di conferma che persiste i dati validati. Il sistema non deve mai interpretare in modo silenzioso un valore malformato o mancante come zero.

\subsection{Requisiti funzionali di configurazione dello scenario}

\RF{Configurazione di uno scenario di ottimizzazione.} \label{rf:config_scenario}
Il sistema deve consentire all'operatore tecnico di costruire uno scenario di ottimizzazione specificando almeno: budget complessivo disponibile, obiettivo da perseguire tra quelli supportati (risparmio energetico, risparmio economico, riduzione di CO\textsubscript{2}) e insieme delle opzioni di intervento candidate. Il sistema deve rendere sempre visibili all'operatore i parametri correnti dello scenario.

\RF{Definizione dei vincoli dello scenario.} \label{rf:vincoli_scenario}
Il sistema deve consentire all'operatore di specificare sullo scenario i vincoli applicabili, tra cui:
\begin{itemize}
    \item coppie di opzioni mutuamente incompatibili;
    \item relazioni di prerequisito tra opzioni (un'opzione richiede che un'altra sia selezionata);
    \item un numero massimo di opzioni selezionabili all'interno di un sottoinsieme specificato;
    \item opzioni obbligatoriamente incluse oppure obbligatoriamente escluse dal portafoglio.
\end{itemize}
Il sistema deve impedire la costruzione di scenari logicamente incoerenti (es. la stessa opzione contemporaneamente inclusa ed esclusa).

\RF{Selezione del sottoinsieme di candidati.} \label{rf:candidati}
Il sistema deve consentire di selezionare in modo puntuale quali opzioni di intervento entrano a far parte dell'insieme dei candidati dello scenario, a partire dall'intero catalogo. Il sistema deve fornire funzioni di selezione e deselezione rapida (es. seleziona tutti, deseleziona tutti) e deve mostrare in modo chiaro il numero di candidati correntemente selezionati.

\subsection{Requisiti funzionali di ottimizzazione e spiegazione}

\RF{Esecuzione dell'ottimizzazione.} \label{rf:ottimizzazione}
Il sistema deve, su richiesta dell'operatore, calcolare il portafoglio ottimo di opzioni di intervento rispetto all'obiettivo dello scenario, nel rispetto di tutti i vincoli specificati. Il sistema deve comunicare in modo esplicito lo stato del calcolo (ottimo dimostrato, ammissibile senza prova di ottimalità, inammissibile, indeterminato per timeout) e, in caso di ottimo, deve produrre la lista degli interventi selezionati corredata degli indicatori aggregati dello scenario (investimento totale, risparmio energetico totale, risparmio economico totale, CO\textsubscript{2} evitata). Il calcolo deve essere deterministico: a parità di scenario, lo stesso portafoglio deve essere prodotto.

\RF{Visualizzazione dei risultati dell'ottimizzazione.} \label{rf:risultati}
Il sistema deve presentare i risultati dell'ottimizzazione con indicatori sintetici dello scenario e con l'elenco dettagliato degli interventi selezionati e non selezionati. Per ciascun intervento deve essere reso visibile almeno il contributo all'obiettivo, il costo, il risparmio energetico, il risparmio economico e la riduzione di CO\textsubscript{2}.

\RF{Spiegazione deterministica delle scelte.} \label{rf:spiegazioni}
Il sistema deve produrre, per ciascuna opzione candidata, una spiegazione deterministica e verificabile della sua selezione o esclusione, classificata in una delle seguenti categorie:
\begin{itemize}
    \item \emph{selezionato} per contributo diretto all'obiettivo;
    \item \emph{selezionato in quanto prerequisito} di altre opzioni selezionate;
    \item \emph{escluso per incompatibilità} con un'opzione selezionata;
    \item \emph{escluso per prerequisito mancante};
    \item \emph{escluso per budget totale insufficiente};
    \item \emph{escluso per budget residuo insufficiente};
    \item \emph{escluso per vincolo di numero massimo};
    \item \emph{escluso per vincolo esplicito di esclusione};
    \item \emph{escluso in quanto non ottimale} (causa specifica non deducibile localmente).
\end{itemize}
La spiegazione deve distinguere esplicitamente i casi in cui la causa è dimostrabile localmente da quelli in cui è classificata come ``non ottimale''.

\RF{Consultazione dello storico delle ottimizzazioni.} \label{rf:storico}
Il sistema deve registrare le esecuzioni di ottimizzazione in uno storico consultabile, mantenendo per ciascuna esecuzione almeno: istante di esecuzione, parametri dello scenario, stato del calcolo, indicatori aggregati e hash canonico del problema.

\subsection{Requisiti funzionali di analisi comparativa}

\RF{Gestione degli scenari salvati.} \label{rf:scenari_salvati}
Il sistema deve consentire all'operatore di salvare lo scenario corrente in due slot di confronto (A e B) e di richiamarli per successive modifiche o esecuzioni. Il sistema deve rendere visibile in ogni momento se ciascuno slot è popolato o vuoto.

\RF{Confronto tra due scenari.} \label{rf:confronto}
Il sistema deve consentire di confrontare due scenari (A e B) e di mostrare, per ciascuno, i parametri principali e gli indicatori aggregati, mettendo in evidenza le differenze (scostamenti di investimento, risparmio energetico, risparmio economico, CO\textsubscript{2}, numero di interventi selezionati).

\RF{Analisi di sensibilità sul budget.} \label{rf:sensibilita}
Il sistema deve consentire di eseguire, a partire dallo scenario corrente, un'analisi di sensibilità sul parametro budget, mostrando come variano l'obiettivo e la composizione del portafoglio al variare del budget disponibile. L'analisi deve classificare ogni opzione in \emph{sempre selezionata}, \emph{mai selezionata} e \emph{selezionata solo per alcune soglie di budget}, e deve essere deterministica.

\subsection{Requisiti funzionali di assistenza AI (opzionali)}

\RF{Proposta assistita di configurazione dello scenario.} \label{rf:ai_proposta}
Il sistema deve offrire, in modo opzionale, la possibilità di ottenere una proposta di configurazione dello scenario a partire da una richiesta in linguaggio naturale dell'operatore (es. ``\emph{vorrei concentrarmi sulla riduzione di CO\textsubscript{2} con 500 mila euro}''). La proposta deve riguardare esclusivamente i parametri di configurazione (budget, obiettivo, insieme dei candidati) e deve essere resa all'operatore corredata da informazioni sull'origine (motore a regole oppure modello linguistico esterno) e da eventuali segnalazioni di incompletezza.

\RF{Approvazione umana della proposta AI.} \label{rf:ai_approvazione}
Il sistema deve richiedere sempre un'atto esplicito di approvazione da parte dell'operatore prima di applicare una proposta AI allo scenario corrente. Il sistema deve consentire all'operatore di accettare, modificare oppure rifiutare la proposta; in nessun caso la proposta AI deve essere applicata automaticamente.

\subsection{Requisiti funzionali di presentazione}

\RF{Cambio del tema di visualizzazione.} \label{rf:tema}
Il sistema deve consentire all'utente di alternare tra un tema chiaro e un tema scuro per l'interfaccia, mantenendo la scelta tra una sessione e la successiva. La preferenza del sistema operativo deve essere utilizzata come valore di default alla prima apertura.
\section{Requisiti Non Funzionali}

\RNF{Isolamento del dominio applicativo.} \label{rnf:isolamento}
Il sistema deve mantenere una separazione netta tra i seguenti livelli: (i) dati di catalogo, (ii) parametri di scenario, (iii) vincoli formali, (iv) motore di ottimizzazione, (v) generazione di spiegazioni, (vi) componente di assistenza AI, (vii) decisione finale dell'utente. Nessun livello deve poter modificare direttamente lo stato dei livelli che non gli competono.

\textbf{Motivazione:} la separazione è un vincolo di dominio del progetto ed è funzionale alla verificabilità del risultato: consente di auditare separatamente dati, calcolo e spiegazioni, e impedisce che la componente AI possa influenzare il risultato matematico.

\RNF{Determinismo e riproducibilità dell'ottimizzazione.} \label{rnf:determinismo}
A parità di parametri di scenario (budget, obiettivo, insieme dei candidati, vincoli, configurazione dei timeout), l'ottimizzazione deve produrre lo stesso portafoglio e gli stessi indicatori aggregati, indipendentemente dall'ordine di presentazione dei candidati. Il sistema deve associare a ogni problema di ottimizzazione un \emph{hash canonico} indipendente dall'ordine degli input.

\textbf{Motivazione:} la riproducibilità è condizione necessaria perché un risultato possa essere discusso in sede decisionale e perché eventuali anomalie possano essere riprodotte e diagnosticate.

\RNF{Prestazioni dell'ottimizzazione.} \label{rnf:prestazioni}
L'ottimizzazione di uno scenario contenente fino a 30 opzioni candidate deve concludersi entro 5 secondi nel 95\% delle esecuzioni, in condizioni di deployment standard (Docker Compose, hardware di riferimento del PoC). In caso di superamento del tempo massimo, il sistema deve comunque restituire uno stato esplicito di ``indeterminato per timeout'', senza bloccare l'interfaccia.

\textbf{Motivazione:} la soglia di 5 secondi è stata concordata per garantire che la funzionalità di analisi what-if possa essere utilizzata in modo interattivo durante riunioni tecniche, senza interrompere il flusso di discussione dell'operatore.

\RNF{Tracciabilità delle decisioni.} \label{rnf:tracciabilita}
Ogni esecuzione di ottimizzazione deve essere persistita ai fini di audit, con: istante di esecuzione, parametri di input, stato del calcolo, indici aggregati, spiegazioni prodotte e hash canonico del problema. Lo storico deve essere consultabile e non modificabile dall'interfaccia utente.

\textbf{Motivazione:} l'auditabilità delle scelte è richiesta per garantire che le decisioni assunte sulla base di TEPO possano essere ricostruite e giustificate in sede di controllo amministrativo, anche a distanza di tempo.

\RNF{Interpretabilità delle spiegazioni.} \label{rnf:interpretabilita}
Ogni opzione candidata deve essere corredata da una spiegazione deterministica riconducibile a una categoria esplicita (selezione diretta, selezione come prerequisito, esclusione per uno specifico vincolo, esclusione non localmente dimostrabile). Il sistema non deve produrre spiegazioni in linguaggio naturale generato in modo non deterministico per classificare la causa di selezione o esclusione.

\textbf{Motivazione:} la trasparenza decisionale è uno degli obiettivi primari dichiarati dal committente; spiegazioni opache o non riproducibili non sarebbero accettabili in sede di presentazione dei risultati alla cittadinanza.

\RNF{Isolamento e opzionalità della componente AI.} \label{rnf:ai_isolamento}
La componente di assistenza AI deve essere confinata alla generazione di proposte di configurazione e non deve avere accesso in scrittura ai dati di catalogo, ai vincoli, ai parametri di scenario applicati, al risultato dell'ottimizzazione, né alle spiegazioni. L'applicativo deve essere pienamente funzionante in assenza di configurazione di un provider LLM e non deve richiedere alcuna chiave API per l'esecuzione della demo locale.

\textbf{Motivazione:} l'app funziona completamente senza API key per scelta architetturale esplicita, al fine di garantire la riproducibilità della demo in contesti offline o in assenza di credenziali esterne, e per evitare che l'attivazione di un LLM possa introdurre comportamenti non verificabili nel perimetro decisionale.

\RNF{Compatibilità.} \label{rnf:compatibilita}
Il sistema deve essere utilizzabile da browser moderni: Chrome versione 100 o superiore, Firefox versione 100 o superiore, Safari versione 15 o superiore, Microsoft Edge versione 100 o superiore. Il design deve essere responsivo e garantire la corretta fruizione dei contenuti su viewport con larghezza compresa tra 320px (mobile) e 1920px (desktop).

\textbf{Motivazione:} le versioni minime indicate corrispondono alle versioni dei browser installate sulle postazioni di lavoro tipicamente utilizzate negli uffici tecnici del Comune e sui dispositivi personali degli operatori, coerentemente con la compatibilità dichiarata per il progetto.

\RNF{Accessibilità.} \label{rnf:accessibilita}
Il sistema deve essere conforme alle linee guida WCAG 2.1 livello AA per quanto riguarda contrasto cromatico, navigazione da tastiera, focus visibile, testi alternativi per le immagini e struttura semantica delle pagine.

\textbf{Motivazione:} la platea degli operatori tecnici include utenti con livelli di familiarità digitale diversificati; il rispetto delle linee guida WCAG 2.1 AA è stato assunto come riferimento perché è lo standard richiesto per i servizi pubblici digitali.

\RNF{Manutenibilità e documentazione delle API.} \label{rnf:manutenibilita}
Il sistema deve esporre le proprie API secondo una descrizione formale consultabile in modalità interattiva (interfaccia Swagger UI) e deve essere organizzato in moduli indipendenti e testabili. Ogni modifica alle API pubbliche deve essere riflessa nella documentazione tecnica.

\textbf{Motivazione:} la manutenibilità del sistema nel medio periodo è un requisito esplicito, in quanto il PoC è destinato a evolvere verso l'integrazione con dati reali del Comune e sarà oggetto di valutazione in sede accademica.

\RNF{Snapshot immutabile delle revisioni dati.} \label{rnf:snapshot}
Ogni esecuzione di ottimizzazione registrata ai fini di audit deve essere ricostruibile a partire dall'hash canonico del problema e dalla revisione del catalogo dati utilizzata al momento dell'esecuzione. Il sistema deve rendere possibile associare a ciascuna esecuzione l'identificativo della revisione dei dati.

\textbf{Motivazione:} senza un riferimento esplicito alla revisione del catalogo, un risultato prodotto in una sessione precedente non sarebbe riproducibile qualora il catalogo fosse successivamente modificato, compromettendo l'affidabilità dell'audit.
\section{Use Case Diagram}

\subsection*{Attori}
\begin{itemize}
    \item \textbf{Operatore Tecnico:} tecnico del Comune incaricato di costruire scenari, eseguire ottimizzazioni e analizzare i risultati.
    \item \textbf{Decisore Comunale:} soggetto (amministratore, dirigente) che consulta i risultati e le spiegazioni a supporto della decisione finale. Nell'attuale implementazione il ruolo coincide con quello dell'Operatore Tecnico.
    \item \textbf{Assistente AI} (opzionale): componente che produce proposte di configurazione a partire da richieste in linguaggio naturale. È modellato come \emph{attore secondario}.
    \item \textbf{Motore di Ottimizzazione} (CP-SAT): sistema esterno che risolve il problema di ottimizzazione combinatoria. Modellato come \emph{attore secondario}.
    \item \textbf{Servizio Mappa:} sistema esterno che fornisce le tile cartografiche utilizzate per la visualizzazione geografica degli edifici.
\end{itemize}

\subsection*{Diagramma principale}

\begin{lstlisting}[language=,caption={Use Case Diagram principale di TEPO}]
@startuml
left to right direction
skinparam packageStyle rectangle

actor "Operatore Tecnico" as OT
actor "Decisore Comunale" as DC
actor "Assistente AI" as AI <<optional>>
actor "Motore di Ottimizzazione\n(CP-SAT)" as SOLVER <<system>>
actor "Servizio Mappa" as MAP <<system>>

rectangle TEPO {
  usecase "UC-01\nConsultare il catalogo\ndegli edifici" as UC01
  usecase "UC-02\nConsultare il catalogo\ndegli interventi" as UC02
  usecase "UC-03\nConfigurare uno\nscenario" as UC03
  usecase "UC-04\nDefinire i vincoli\ndello scenario" as UC04
  usecase "UC-05\nEseguire\nl'ottimizzazione" as UC05
  usecase "UC-06\nConsultare risultati\ne spiegazioni" as UC06
  usecase "UC-07\nSalvare e confrontare\nscenari A/B" as UC07
  usecase "UC-08\nAnalisi di\nsensibilita sul budget" as UC08
  usecase "UC-09\nOttenere una proposta\nassistita dall'AI" as UC09
  usecase "UC-10\nApprovare o rifiutare\nla proposta AI" as UC10
  usecase "UC-11\nImportare dati\nda CSV" as UC11
  usecase "UC-12\nConsultare lo storico\ndelle ottimizzazioni" as UC12
  usecase "UC-13\nVisualizzare gli edifici\nsu mappa" as UC13
}

OT --> UC01
OT --> UC02
OT --> UC03
OT --> UC04
OT --> UC05
OT --> UC06
OT --> UC07
OT --> UC08
OT --> UC09
OT --> UC10
OT --> UC11
OT --> UC12
OT --> UC13
DC --> UC06
DC --> UC07

UC05 ..> SOLVER : <<include>>
UC13 ..> MAP : <<include>>
UC09 ..> AI : <<include>>
UC10 .> UC09 : <<extend>>
UC07 ..> UC06 : <<include>>
UC08 ..> UC05 : <<include>>

@enduml
\end{lstlisting}

% IMMAGINE: img/UseCaseDiagram.png -- Diagramma dei casi d'uso principale di TEPO

\subsection{Descrizione dei casi d'uso principali}

\paragraph{\large UC-03: Configurare uno scenario di ottimizzazione}

\paragraph{RF collegati:} RF\ref{rf:config_scenario}, RF\ref{rf:candidati}, RF\ref{rf:vincoli_scenario}.

\paragraph{Attore primario:} Operatore Tecnico.

\paragraph{Obiettivo:} definire in modo completo i parametri (budget, obiettivo, candidati, vincoli) di uno scenario pronto per l'esecuzione dell'ottimizzazione.

\paragraph{Precondizioni:} il catalogo degli edifici e degli interventi è popolato; l'utente ha accesso all'interfaccia di configurazione dello scenario.

\paragraph{Postcondizioni:} lo scenario corrente contiene un budget valido, un obiettivo selezionato, un insieme non vuoto di candidati e un insieme di vincoli ben formati.

\paragraph{Flusso principale:}
\begin{enumerate}
    \item L'utente apre la pagina di configurazione dello scenario.
    \item Il sistema mostra i parametri correnti dello scenario (budget, obiettivo, candidati, vincoli).
    \item L'utente inserisce o modifica il budget complessivo.
    \item L'utente seleziona l'obiettivo tra risparmio energetico, risparmio economico e riduzione di CO\textsubscript{2}.
    \item L'utente seleziona le opzioni di intervento candidate a partire dal catalogo. [Estensione 1]
    \item L'utente aggiunge eventuali vincoli di incompatibilità, prerequisito, numero massimo, inclusione o esclusione. [Estensione 2]
    \item Il sistema verifica la coerenza logica dei parametri. [Eccezione 1]
    \item Il sistema conferma all'utente che lo scenario è pronto per l'ottimizzazione.
\end{enumerate}

\paragraph{Estensioni:}
\begin{enumerate}
    \item L'utente può utilizzare le funzioni di selezione rapida (seleziona tutti / deseleziona tutti) oppure filtrare il catalogo per edificio e tipologia prima di selezionare i candidati.
    \item L'utente può rimuovere un vincolo precedentemente aggiunto.
\end{enumerate}

\paragraph{Eccezioni:}
\begin{enumerate}
    \item Se il budget è nullo o negativo, se non è stato selezionato alcun obiettivo, oppure se la stessa opzione risulta contemporaneamente inclusa ed esclusa, il sistema blocca il salvataggio e mostra un messaggio di errore specifico.
\end{enumerate}

\paragraph{\large UC-05: Eseguire l'ottimizzazione e consultare le spiegazioni}

\paragraph{RF collegati:} RF\ref{rf:ottimizzazione}, RF\ref{rf:risultati}, RF\ref{rf:spiegazioni}.

\paragraph{Attore primario:} Operatore Tecnico.

\paragraph{Attori secondari:} Motore di Ottimizzazione.

\paragraph{Obiettivo:} ottenere il portafoglio ottimo di interventi per lo scenario corrente e comprenderne le motivazioni.

\paragraph{Precondizioni:} lo scenario corrente è completo e coerente (UC-03 concluso con successo).

\paragraph{Postcondizioni:} il sistema mostra il portafoglio ottimo, gli indicatori aggregati e le spiegazioni per ciascuna opzione candidata. L'esecuzione è registrata nello storico delle ottimizzazioni.

\paragraph{Flusso principale:}
\begin{enumerate}
    \item L'utente richiede l'esecuzione dell'ottimizzazione.
    \item Il sistema valida nuovamente lo scenario.
    \item Il sistema invoca il motore di ottimizzazione. [Estensione 1]
    \item Il motore di ottimizzazione restituisce una soluzione con lo stato del calcolo.
    \item Il sistema calcola le spiegazioni deterministiche per ciascuna opzione candidata.
    \item Il sistema registra l'esecuzione nello storico ai fini di audit.
    \item Il sistema presenta all'utente il portafoglio ottimo, gli indicatori aggregati e le spiegazioni.
\end{enumerate}

\paragraph{Estensioni:}
\begin{enumerate}
    \item In caso di più soluzioni ottime equivalenti, il sistema applica criteri di tie-break deterministici (minimizzazione del costo, minimizzazione del numero di interventi, preferenza per identificativi crescenti) e ne dà evidenza all'utente.
\end{enumerate}

\paragraph{Eccezioni:}
\begin{enumerate}
    \item Se il problema risulta inammissibile, il sistema mostra lo stato ``inammissibile'' e propone all'utente una diagnostica sulle possibili cause (es. vincolo di inclusione su un'opzione non candidata).
    \item Se il calcolo supera il tempo massimo, il sistema mostra lo stato ``indeterminato per timeout'' e non presenta alcun portafoglio come definitivo.
\end{enumerate}

\paragraph{\large UC-07: Salvare e confrontare due scenari}

\paragraph{RF collegati:} RF\ref{rf:scenari_salvati}, RF\ref{rf:confronto}.

\paragraph{Attore primario:} Operatore Tecnico.

\paragraph{Attore secondario:} Decisore Comunale.

\paragraph{Obiettivo:} confrontare in modo strutturato due configurazioni alternative (scenario A e scenario B) per orientare la decisione.

\paragraph{Precondizioni:} almeno due scenari sono stati configurati ed eventualmente ottimizzati.

\paragraph{Postcondizioni:} il sistema mostra in modo comparato i parametri e gli indicatori aggregati dei due scenari, evidenziando gli scostamenti.

\paragraph{Flusso principale:}
\begin{enumerate}
    \item L'utente salva lo scenario corrente nello slot A.
    \item L'utente modifica lo scenario corrente e lo salva nello slot B. [Estensione 1]
    \item L'utente accede alla pagina di confronto.
    \item Il sistema recupera dagli slot A e B i parametri e gli indicatori aggregati delle rispettive esecuzioni.
    \item Il sistema mostra i due scenari in forma tabellare, evidenziando gli scostamenti.
\end{enumerate}

\paragraph{Estensioni:}
\begin{enumerate}
    \item L'utente può sovrascrivere lo slot A o lo slot B con lo scenario corrente in qualunque momento, dopo conferma esplicita.
\end{enumerate}

\paragraph{\large UC-09: Ottenere una proposta assistita dall'AI}

\paragraph{RF collegati:} RF\ref{rf:ai_proposta}, RF\ref{rf:ai_approvazione}.

\paragraph{Attore primario:} Operatore Tecnico.

\paragraph{Attore secondario:} Assistente AI (opzionale).

\paragraph{Obiettivo:} ottenere una proposta preliminare di configurazione dello scenario a partire da una richiesta in linguaggio naturale.

\paragraph{Precondizioni:} l'utente si trova nella pagina di configurazione dello scenario.

\paragraph{Postcondizioni:} il sistema mostra una proposta di configurazione corredata da origine, eventuali segnalazioni e richiesta esplicita di approvazione umana.

\paragraph{Flusso principale:}
\begin{enumerate}
    \item L'utente apre il pannello di assistenza AI.
    \item L'utente digita la propria richiesta in linguaggio naturale.
    \item Il sistema inoltra la richiesta al componente di assistenza (motore a regole oppure, se configurato, modello linguistico esterno). [Estensione 1]
    \item Il componente di assistenza restituisce una proposta di configurazione con origine, eventuali informazioni mancanti ed eventuali avvisi.
    \item Il sistema mostra la proposta all'utente senza applicarla.
    \item L'utente accetta, modifica oppure rifiuta la proposta. [Eccezione 1]
    \item In caso di accettazione, il sistema applica i parametri proposti allo scenario corrente e ne dà evidenza all'utente.
\end{enumerate}

\paragraph{Estensioni:}
\begin{enumerate}
    \item Se il modello linguistico esterno non è raggiungibile o restituisce una risposta non conforme, il sistema ricade sul motore a regole e segnala l'evento all'utente.
\end{enumerate}

\paragraph{Eccezioni:}
\begin{enumerate}
    \item Se l'utente rifiuta la proposta, il sistema non modifica lo scenario corrente.
\end{enumerate}
\section{User Stories}

\textbf{Premessa:} in fase di implementazione alcune funzionalità potrebbero essere semplificate. Eventuali differenze tra progettazione e implementazione verranno chiarite nel documento D2.

\US{Consultazione del catalogo degli edifici} \label{us:edifici}
Come \textit{operatore tecnico del Comune} voglio consultare il catalogo degli edifici pubblici per individuare rapidamente gli immobili su cui sono già stati identificati interventi candidati. \\ (Riferito a RF\ref{rf:catalogo_edifici})

\paragraph{Criteri di accettazione:}
\begin{itemize}
    \item \textbf{AC-US01-01:} Dato che l'utente accede alla pagina ``Edifici'', quando la pagina è caricata, allora il sistema mostra l'elenco completo degli edifici con identificativo, denominazione, indirizzo, tipologia d'uso e numero di interventi associati.
    \item \textbf{AC-US01-02:} Dato che l'utente è sulla pagina ``Edifici'', quando applica un filtro per tipologia d'uso, allora il sistema mostra solo gli edifici corrispondenti.
    \item \textbf{AC-US01-03:} Dato che l'utente è sulla pagina ``Edifici'', quando clicca sull'intestazione di una colonna, allora il sistema ordina l'elenco per la colonna selezionata.
\end{itemize}

\US{Visualizzazione geografica degli edifici} \label{us:mappa}
Come \textit{operatore tecnico del Comune} voglio visualizzare gli edifici su una mappa per comprenderne la distribuzione territoriale e valutare interventi in aree omogenee. \\ (Riferito a RF\ref{rf:mappa})

\paragraph{Criteri di accettazione:}
\begin{itemize}
    \item \textbf{AC-US02-01:} Dato che l'utente accede alla pagina ``Mappa'', quando la pagina è caricata, allora il sistema mostra gli edifici come marker sulla mappa, posizionati alle coordinate registrate.
    \item \textbf{AC-US02-02:} Dato che l'utente visualizza la mappa con uno scenario caricato, quando i marker sono renderizzati, allora il sistema colora i marker in base allo stato dell'edificio nello scenario (con interventi selezionati, esclusi, non valutati).
\end{itemize}

\US{Consultazione e ricerca degli interventi} \label{us:interventi}
Come \textit{operatore tecnico del Comune} voglio consultare e filtrare le opzioni di intervento disponibili per costruire un sottoinsieme di candidati ragionato. \\ (Riferito a RF\ref{rf:catalogo_interventi}, RF\ref{rf:filtri_interventi})

\paragraph{Criteri di accettazione:}
\begin{itemize}
    \item \textbf{AC-US03-01:} Dato che l'utente accede alla pagina ``Interventi'', quando la pagina è caricata, allora il sistema mostra l'elenco completo degli interventi con edificio, tipologia, costo, risparmio energetico, risparmio economico, CO\textsubscript{2} e payback.
    \item \textbf{AC-US03-02:} Dato che l'utente è sulla pagina ``Interventi'', quando applica uno o più filtri (ricerca testuale, edificio, tipologia, intervallo di costo, stato), allora il sistema mostra solo gli interventi che soddisfano tutti i criteri applicati.
    \item \textbf{AC-US03-03:} Dato che l'utente è sulla pagina ``Interventi'', quando la ricerca non produce risultati, allora il sistema mostra un messaggio informativo.
\end{itemize}

\US{Importazione di dati di catalogo da CSV} \label{us:import}
Come \textit{operatore tecnico del Comune} voglio poter caricare dati di catalogo da file CSV per non dover inserire manualmente gli edifici e gli interventi. \\ (Riferito a RF\ref{rf:import_csv})

\paragraph{Criteri di accettazione:}
\begin{itemize}
    \item \textbf{AC-US04-01:} Dato che l'utente accede alla pagina ``Importazione'', quando carica un file CSV, allora il sistema mostra un'anteprima delle righe valide, delle righe scartate e delle relative motivazioni.
    \item \textbf{AC-US04-02:} Dato che il sistema ha prodotto un'anteprima di validazione, quando l'utente conferma l'importazione, allora il sistema persiste solo le righe valide.
    \item \textbf{AC-US04-03:} Dato che una riga contiene un valore numerico malformato, quando il sistema esegue la validazione, allora la riga viene scartata con motivazione esplicita e non viene interpretata come zero.
\end{itemize}

\US{Configurazione di uno scenario di ottimizzazione} \label{us:scenario}
Come \textit{operatore tecnico del Comune} voglio configurare uno scenario di ottimizzazione specificando budget, obiettivo e candidati per ottenere un portafoglio ottimo di interventi. \\ (Riferito a RF\ref{rf:config_scenario}, RF\ref{rf:candidati})

\paragraph{Criteri di accettazione:}
\begin{itemize}
    \item \textbf{AC-US05-01:} Dato che l'utente è nella pagina ``Scenario'', quando modifica il budget, allora il sistema aggiorna il valore visualizzato e lo mantiene in memoria fino al salvataggio.
    \item \textbf{AC-US05-02:} Dato che l'utente è nella pagina ``Scenario'', quando seleziona un obiettivo tra risparmio energetico, risparmio economico e riduzione di CO\textsubscript{2}, allora il sistema evidenzia l'obiettivo attivo.
    \item \textbf{AC-US05-03:} Dato che l'utente è nella pagina ``Scenario'', quando seleziona o deseleziona candidati, allora il sistema aggiorna il contatore dei candidati selezionati.
\end{itemize}

\US{Definizione dei vincoli dello scenario} \label{us:vincoli}
Come \textit{operatore tecnico del Comune} voglio specificare i vincoli dello scenario (incompatibilità, prerequisiti, numero massimo, inclusione ed esclusione) per riflettere i vincoli tecnici e di policy dell'Amministrazione. \\ (Riferito a RF\ref{rf:vincoli_scenario})

\paragraph{Criteri di accettazione:}
\begin{itemize}
    \item \textbf{AC-US06-01:} Dato che l'utente è nella pagina ``Scenario'', quando aggiunge una coppia di opzioni incompatibili, allora il sistema memorizza il vincolo e lo mostra nella lista dei vincoli correnti.
    \item \textbf{AC-US06-02:} Dato che l'utente è nella pagina ``Scenario'', quando specifica che un'opzione B richiede un'opzione A come prerequisito, allora il sistema memorizza il vincolo.
    \item \textbf{AC-US06-03:} Dato che l'utente tenta di salvare uno scenario con la stessa opzione contemporaneamente inclusa ed esclusa, allora il sistema blocca il salvataggio e mostra un messaggio di errore.
\end{itemize}

\US{Esecuzione dell'ottimizzazione} \label{us:ottimizza}
Come \textit{operatore tecnico del Comune} voglio eseguire l'ottimizzazione dello scenario per ottenere il portafoglio ottimo rispetto all'obiettivo scelto. \\ (Riferito a RF\ref{rf:ottimizzazione}, RF\ref{rf:risultati})

\paragraph{Criteri di accettazione:}
\begin{itemize}
    \item \textbf{AC-US07-01:} Dato che lo scenario corrente è valido, quando l'utente richiede l'ottimizzazione, allora il sistema presenta il portafoglio selezionato, gli indicatori aggregati e lo stato del calcolo.
    \item \textbf{AC-US07-02:} Dato che lo scenario corrente contiene opzioni candidate, quando l'utente esegue due volte l'ottimizzazione senza modificare i parametri, allora il sistema restituisce lo stesso portafoglio e lo stesso hash canonico del problema.
    \item \textbf{AC-US07-03:} Dato che il calcolo supera il tempo massimo configurato, quando il sistema produce il risultato, allora lo stato riportato è ``indeterminato per timeout'' e l'interfaccia non resta bloccata.
\end{itemize}

\US{Consultazione delle spiegazioni} \label{us:spiegazioni}
Come \textit{operatore tecnico del Comune} voglio leggere la motivazione della selezione o dell'esclusione di ciascun intervento per poter giustificare la scelta in sede amministrativa. \\ (Riferito a RF\ref{rf:spiegazioni})

\paragraph{Criteri di accettazione:}
\begin{itemize}
    \item \textbf{AC-US08-01:} Dato che l'utente visualizza i risultati dell'ottimizzazione, quando scorre la sezione ``Spiegazioni'', allora il sistema mostra per ciascuna opzione candidata la categoria di motivazione.
    \item \textbf{AC-US08-02:} Dato che un'opzione è stata esclusa per budget insufficiente, quando l'utente consulta la spiegazione, allora la categoria mostrata è ``esclusa per budget totale insufficiente'' oppure ``esclusa per budget residuo insufficiente''.
    \item \textbf{AC-US08-03:} Dato che la causa dell'esclusione non è localmente dimostrabile, quando l'utente consulta la spiegazione, allora la categoria mostrata è ``esclusa in quanto non ottimale''.
\end{itemize}

\US{Salvataggio di due scenari A/B} \label{us:salva_ab}
Come \textit{operatore tecnico del Comune} voglio salvare lo scenario corrente in due slot di confronto per poter costruire e confrontare alternative di investimento. \\ (Riferito a RF\ref{rf:scenari_salvati})

\paragraph{Criteri di accettazione:}
\begin{itemize}
    \item \textbf{AC-US09-01:} Dato che l'utente ha configurato uno scenario, quando lo salva nello slot A, allora il sistema memorizza i parametri correnti nello slot A e lo indica visivamente.
    \item \textbf{AC-US09-02:} Dato che lo slot A è popolato, quando l'utente salva un nuovo scenario nello slot B, allora il sistema mantiene invariato il contenuto dello slot A.
\end{itemize}

\US{Confronto tra due scenari} \label{us:confronto}
Come \textit{decisore comunale} voglio confrontare due scenari A e B per valutare alternative di allocazione del budget con indicatori omogenei. \\ (Riferito a RF\ref{rf:confronto})

\paragraph{Criteri di accettazione:}
\begin{itemize}
    \item \textbf{AC-US10-01:} Dato che entrambi gli slot A e B sono popolati, quando l'utente accede alla pagina ``Confronto'', allora il sistema mostra i parametri e gli indicatori aggregati di entrambi.
    \item \textbf{AC-US10-02:} Dato che entrambi gli scenari sono visualizzati, quando il sistema mostra gli indicatori, allora evidenzia gli scostamenti (\emph{delta}) tra A e B per investimento, risparmio energetico, risparmio economico, CO\textsubscript{2} e numero di interventi.
\end{itemize}

\US{Analisi di sensibilità sul budget} \label{us:sensibilita}
Come \textit{operatore tecnico del Comune} voglio analizzare come cambia il portafoglio ottimo al variare del budget per individuare la soglia oltre la quale il rendimento marginale decresce. \\ (Riferito a RF\ref{rf:sensibilita})

\paragraph{Criteri di accettazione:}
\begin{itemize}
    \item \textbf{AC-US11-01:} Dato che lo scenario corrente è valido, quando l'utente richiede l'analisi di sensibilità, allora il sistema esegue una griglia deterministica di ottimizzazioni su valori di budget crescenti.
    \item \textbf{AC-US11-02:} Dato che l'analisi è conclusa, quando l'utente visualizza i risultati, allora il sistema classifica ciascuna opzione in ``sempre selezionata'', ``mai selezionata'' e ``selezionata solo per alcune soglie di budget''.
\end{itemize}

\US{Proposta assistita dall'AI} \label{us:ai}
Come \textit{operatore tecnico del Comune} voglio descrivere a parole l'obiettivo che ho in mente per ottenere una proposta preliminare di configurazione dello scenario, così da ridurre il tempo necessario a impostarla. \\ (Riferito a RF\ref{rf:ai_proposta})

\paragraph{Criteri di accettazione:}
\begin{itemize}
    \item \textbf{AC-US12-01:} Dato che l'utente digita una richiesta in linguaggio naturale, quando invia la richiesta, allora il sistema restituisce una proposta con budget, obiettivo e insieme di candidati, corredata da origine (regole oppure modello esterno).
    \item \textbf{AC-US12-02:} Dato che il componente esterno non è configurato, quando l'utente invia la richiesta, allora il sistema produce comunque una proposta attraverso il motore a regole e ne indica esplicitamente l'origine.
\end{itemize}

\US{Approvazione umana della proposta AI} \label{us:ai_appr}
Come \textit{operatore tecnico del Comune} voglio poter accettare, modificare o rifiutare esplicitamente una proposta AI prima che questa venga applicata allo scenario, per mantenere il pieno controllo sulle decisioni. \\ (Riferito a RF\ref{rf:ai_approvazione})

\paragraph{Criteri di accettazione:}
\begin{itemize}
    \item \textbf{AC-US13-01:} Dato che il sistema ha prodotto una proposta AI, quando l'utente la visualizza, allora la proposta non è applicata automaticamente allo scenario corrente.
    \item \textbf{AC-US13-02:} Dato che l'utente clicca ``Applica'', quando i parametri della proposta sono validi, allora il sistema aggiorna lo scenario corrente e segnala l'avvenuta applicazione.
    \item \textbf{AC-US13-03:} Dato che l'utente clicca ``Rifiuta'', quando la proposta è rifiutata, allora lo scenario corrente resta invariato.
\end{itemize}

\US{Cambio tema di visualizzazione} \label{us:tema}
Come \textit{utente} voglio cambiare tema di visualizzazione tra chiaro e scuro per adattare l'interfaccia alle mie condizioni di utilizzo e preferenze. \\ (Riferito a RF\ref{rf:tema})

\paragraph{Criteri di accettazione:}
\begin{itemize}
    \item \textbf{AC-US14-01:} Dato che l'utente clicca sull'interruttore di tema, quando il tema corrente è quello chiaro, allora il sistema applica il tema scuro all'intera interfaccia.
    \item \textbf{AC-US14-02:} Dato che l'utente ha selezionato un tema, quando chiude e riapre l'applicazione, allora il sistema mantiene il tema selezionato.
\end{itemize}
\section{Design Front-end}

\textbf{Premessa:} in questa sezione vengono riportati dei mockup del frontend dell'applicazione, ossia prototipi di quella che sarà l'interfaccia utente di \textbf{TEPO}. L'effettiva realizzazione del frontend può riportare leggere differenze e maggiori dettagli rispetto agli esempi proposti.

Tutti i mockup rispecchiano il requisito non funzionale di \textbf{Usabilità} e adottano un insieme coerente di token di design (colori, spaziature, tipografia, stati di focus). L'utilizzo del sistema risulta intuitivo perché ogni azione è compiuta premendo pulsanti con indicazioni esplicite sulla loro funzionalità. È disponibile anche un tema scuro, selezionabile dall'utente.

\subsection{Header e navigazione}
\begin{figure}[H]
    \centering
    \includegraphics[width=0.95\linewidth]{img/Header.png}
    \caption{Header di TEPO -- barra di navigazione principale}
    \label{fig:header}
\end{figure}

Nella Figura \ref{fig:header} è riportato un mockup dell'header delle varie pagine dell'applicazione. L'header espone le sezioni principali: \emph{Dashboard}, \emph{Mappa}, \emph{Edifici}, \emph{Interventi}, \emph{Scenario}, \emph{Confronto}, \emph{Importazione}, \emph{Sensibilità}. In alto a destra è presente l'interruttore per il cambio tema (chiaro/scuro) e un badge che segnala in modo permanente che i dati mostrati sono ``DEMO -- dati sintetici''.

% IMMAGINE: img/Header.png -- Header di TEPO con navigazione principale e badge "DEMO - dati sintetici"

\subsection{Dashboard}
\begin{figure}[H]
    \centering
    \includegraphics[width=0.95\linewidth]{img/SchermataDashboard.png}
    \caption{Dashboard di TEPO}
    \label{fig:dashboard}
\end{figure}

Nella Figura \ref{fig:dashboard} è riportato un mockup della dashboard iniziale. Mostra i KPI sintetici del catalogo (numero edifici, numero interventi, tipologie rappresentate) e le azioni rapide verso la costruzione dello scenario. Ogni indicatore è accompagnato dall'unità di misura. La schermata è collegata ai seguenti requisiti:
\begin{itemize}
    \item \textbf{RF\ref{rf:catalogo_edifici}:} consultazione sintetica del catalogo edifici.
    \item \textbf{RF\ref{rf:catalogo_interventi}:} consultazione sintetica del catalogo interventi.
    \item \textbf{RF\ref{rf:config_scenario}:} accesso rapido alla configurazione di uno scenario.
\end{itemize}

Collegamenti principali: RF\ref{rf:catalogo_edifici}, RF\ref{rf:catalogo_interventi}, RF\ref{rf:config_scenario}; UC-01, UC-02, UC-03; US-01, US-03, US-05.

\subsection{Mappa degli edifici}
\begin{figure}[H]
    \centering
    \includegraphics[width=0.95\linewidth]{img/SchermataMappa.png}
    \caption{Vista cartografica degli edifici}
    \label{fig:mappa}
\end{figure}

Nella Figura \ref{fig:mappa} è riportato un mockup della vista cartografica: sulla sinistra la mappa con i marker degli edifici; sulla destra la legenda e la tabella compatta degli edifici con stato corrente rispetto a uno scenario. I colori dei marker distinguono edifici con interventi selezionati, esclusi, non valutati.

Collegamenti principali: RF\ref{rf:mappa}; UC-13; US-02.

\subsection{Catalogo degli edifici}
\begin{figure}[H]
    \centering
    \includegraphics[width=0.95\linewidth]{img/SchermataEdifici.png}
    \caption{Catalogo degli edifici}
    \label{fig:edifici}
\end{figure}

Nella Figura \ref{fig:edifici} è riportato un mockup del catalogo edifici, con filtro per tipologia, ricerca testuale e ordinamento per colonna. La tabella mostra identificativo, denominazione, indirizzo, tipologia, superficie, numero di interventi e provenienza del dato.

Collegamenti principali: RF\ref{rf:catalogo_edifici}; UC-01; US-01.

\subsection{Catalogo degli interventi}
\begin{figure}[H]
    \centering
    \includegraphics[width=0.95\linewidth]{img/SchermataInterventi.png}
    \caption{Catalogo degli interventi}
    \label{fig:interventi}
\end{figure}

Nella Figura \ref{fig:interventi} è riportato un mockup del catalogo interventi con filtri combinabili: ricerca testuale, edificio di riferimento, tipologia, intervallo di costo, stato di valutazione nello scenario corrente. Ogni riga riporta costo, risparmio energetico, risparmio economico, CO\textsubscript{2} e payback semplice.

Collegamenti principali: RF\ref{rf:catalogo_interventi}, RF\ref{rf:filtri_interventi}; UC-02; US-03.

\subsection{Scenario Builder}
\begin{figure}[H]
    \centering
    \includegraphics[width=0.95\linewidth]{img/SchermataScenario.png}
    \caption{Scenario Builder -- configurazione dello scenario}
    \label{fig:scenario}
\end{figure}

Nella Figura \ref{fig:scenario} è riportato un mockup dello \emph{Scenario Builder}. La schermata espone: (i) i campi per budget e obiettivo; (ii) la tabella delle opzioni candidate con possibilità di toggle e selezione rapida; (iii) la sezione dei vincoli applicati (incompatibilità, prerequisiti, numero massimo, inclusione/esclusione); (iv) un pannello di assistenza AI collassabile; (v) il pulsante di esecuzione dell'ottimizzazione.

Collegamenti principali: RF\ref{rf:config_scenario}, RF\ref{rf:candidati}, RF\ref{rf:vincoli_scenario}, RF\ref{rf:ai_proposta}, RF\ref{rf:ai_approvazione}; UC-03, UC-04, UC-09; US-05, US-06, US-12, US-13.

\subsection{Risultati e spiegazioni}
\begin{figure}[H]
    \centering
    \includegraphics[width=0.95\linewidth]{img/SchermataRisultati.png}
    \caption{Pagina dei risultati dell'ottimizzazione}
    \label{fig:risultati}
\end{figure}

Nella Figura \ref{fig:risultati} è riportato un mockup della pagina dei risultati: uno stato del solver ben visibile (badge \emph{OPTIMAL} o altro stato), KPI aggregati (investimento totale, risparmio energetico, risparmio economico, CO\textsubscript{2}, numero di interventi selezionati), tabella degli interventi selezionati e sezione ``Perché questa soluzione?'' con le spiegazioni per ciascuna opzione candidata.

Collegamenti principali: RF\ref{rf:risultati}, RF\ref{rf:spiegazioni}; UC-05, UC-06; US-07, US-08.

\subsection{Confronto tra scenari A/B}
\begin{figure}[H]
    \centering
    \includegraphics[width=0.95\linewidth]{img/SchermataConfronto.png}
    \caption{Pagina di confronto tra due scenari}
    \label{fig:confronto}
\end{figure}

Nella Figura \ref{fig:confronto} è riportato un mockup della pagina di confronto: quattro KPI \emph{delta} (Δ investimento, Δ energia, Δ economia, Δ numero di interventi) e una tabella comparativa con i parametri e gli indicatori dei due scenari. I badge nei header delle sezioni di differenza evidenziano gli interventi aggiunti (verde) o rimossi (arancio) passando da A a B.

Collegamenti principali: RF\ref{rf:scenari_salvati}, RF\ref{rf:confronto}; UC-07; US-09, US-10.

\subsection{Analisi di sensibilità sul budget}
\begin{figure}[H]
    \centering
    \includegraphics[width=0.95\linewidth]{img/SchermataSensibilita.png}
    \caption{Analisi di sensibilità sul budget}
    \label{fig:sensibilita}
\end{figure}

Nella Figura \ref{fig:sensibilita} è riportato un mockup dell'analisi di sensibilità. Sono visibili: il form di avvio (budget minimo, massimo, numero di punti), i KPI sintetici dell'analisi, una tabella con barra proporzionale della frequenza di selezione per ciascuna opzione e tre sezioni dedicate a ``Sempre selezionate'', ``Mai selezionate'' e ``Swing''.

Collegamenti principali: RF\ref{rf:sensibilita}; UC-08; US-11.

\subsection{Importazione dati da CSV}
\begin{figure}[H]
    \centering
    \includegraphics[width=0.95\linewidth]{img/SchermataImport.png}
    \caption{Pagina di importazione dati da CSV}
    \label{fig:import}
\end{figure}

Nella Figura \ref{fig:import} è riportato un mockup della pagina di importazione: descrizione dell'header atteso, pulsante ``Carica esempio'', textarea in monospace per il contenuto CSV, badge con conteggio righe valide e invalide, e pulsante di conferma importazione.

Collegamenti principali: RF\ref{rf:import_csv}; UC-11; US-04.
\section{User Flows}

In questa sezione vengono riportati gli User Flow relativi ai principali comportamenti degli utenti all'interno dell'applicazione. Riportiamo nella seguente figura la legenda dei simboli utilizzati.

% IMMAGINE: img/UF_legenda.png -- Legenda dei simboli usati negli user flow

\begin{figure}[H]
    \centering
    \includegraphics[width=0.7\linewidth]{img/UF_legenda.png}
    \caption{Legenda dei simboli utilizzati negli User Flow}
\end{figure}

\UF{Configurazione scenario, proposta AI, approvazione umana, ottimizzazione e analisi delle spiegazioni} \label{uf:main}
Associato alle user stories US\ref{us:scenario}, US\ref{us:vincoli}, US\ref{us:ottimizza}, US\ref{us:spiegazioni}, US\ref{us:ai}, US\ref{us:ai_appr} e ai requisiti funzionali RF\ref{rf:config_scenario}, RF\ref{rf:vincoli_scenario}, RF\ref{rf:candidati}, RF\ref{rf:ai_proposta}, RF\ref{rf:ai_approvazione}, RF\ref{rf:ottimizzazione}, RF\ref{rf:spiegazioni}.

La Figura \ref{figuf:main} mostra il percorso principale dell'operatore tecnico, che parte dalla configurazione dello scenario e termina con l'analisi delle spiegazioni prodotte dal solver. Il flusso esplicita il passaggio di approvazione umana della proposta AI, che rappresenta un punto architetturalmente vincolante del sistema.

\begin{figure}[H]
    \centering
    \includegraphics[width=1\linewidth]{img/UF_main.png}
    \caption{Configurazione scenario, approvazione AI, ottimizzazione e spiegazioni}
    \label{figuf:main}
\end{figure}

\noindent
Di seguito si riporta la specifica del flusso in forma di diagramma di attività, in modo da chiarire l'ordine delle decisioni e delle transizioni principali.

\begin{lstlisting}[language=,caption={Activity Diagram del flusso principale}]
@startuml
start
:Operatore apre Scenario Builder;
:Configura budget, obiettivo, candidati;
if (Vuole assistenza AI?) then (si)
  :Digita richiesta in linguaggio naturale;
  :Sistema produce proposta AI;
  :Sistema mostra proposta (non applicata);
  if (Operatore approva?) then (si)
    :Sistema applica parametri proposti;
  else (no)
    :Scenario resta invariato;
  endif
else (no)
endif
:Operatore definisce eventuali vincoli;
if (Scenario coerente?) then (si)
  :Operatore richiede ottimizzazione;
  :Sistema invoca il motore di ottimizzazione;
  :Sistema calcola spiegazioni deterministiche;
  :Sistema registra esecuzione in audit;
  :Sistema mostra risultati e spiegazioni;
else (no)
  :Sistema mostra messaggio di errore;
endif
stop
@enduml
\end{lstlisting}

\UF{Confronto tra due scenari alternativi} \label{uf:confronto}
Associato alle user stories US\ref{us:salva_ab} e US\ref{us:confronto} e ai requisiti funzionali RF\ref{rf:scenari_salvati}, RF\ref{rf:confronto}.

La Figura \ref{figuf:confronto} mostra il percorso di costruzione di due alternative (scenario A e scenario B) e il confronto comparato dei risultati. Il decisore comunale è tipicamente fruitore finale di questa vista.

\begin{figure}[H]
    \centering
    \includegraphics[width=1\linewidth]{img/UF_confronto.png}
    \caption{Confronto tra due scenari A/B}
    \label{figuf:confronto}
\end{figure}

\UF{Importazione dati di catalogo da CSV} \label{uf:import}
Associato alla user story US\ref{us:import} e al requisito funzionale RF\ref{rf:import_csv}.

La Figura \ref{figuf:import} mostra il percorso di importazione in due fasi (validazione a secco, poi conferma) con evidenza delle righe scartate e dei motivi di scarto.

\begin{figure}[H]
    \centering
    \includegraphics[width=1\linewidth]{img/UF_import.png}
    \caption{Importazione dati di catalogo da CSV}
    \label{figuf:import}
\end{figure}

\UF{Analisi di sensibilità sul budget} \label{uf:sensibilita}
Associato alla user story US\ref{us:sensibilita} e al requisito funzionale RF\ref{rf:sensibilita}.

La Figura \ref{figuf:sensibilita} mostra il percorso di analisi di sensibilità: definizione della griglia di budget, esecuzione dello sweep, lettura della classificazione di stabilità delle opzioni.

\begin{figure}[H]
    \centering
    \includegraphics[width=1\linewidth]{img/UF_sensibilita.png}
    \caption{Analisi di sensibilità sul budget}
    \label{figuf:sensibilita}
\end{figure}
\section{Traceability}

La matrice di tracciabilità verifica i collegamenti tra i requisiti funzionali e non funzionali del D1 e gli artefatti che li dettagliano (Use Case, User Story, Acceptance Criteria, Mockup). Ogni RF deve essere collegato ad almeno un artefatto; gli RNF sono tracciati verso gli artefatti pertinenti (mockup, user flow o, quando non applicabile, verso la sezione di motivazione).

\begin{longtable}{p{1.2cm} p{2.2cm} p{2.2cm} p{3.4cm} p{2.6cm} p{4cm}}
\caption{Matrice di tracciabilità di TEPO.}\\
\toprule
\textbf{RF/RNF} & \textbf{Use Case} & \textbf{User Story} & \textbf{Acceptance Criteria} & \textbf{Mockup} & \textbf{Note} \\
\midrule
\endfirsthead
\toprule
\textbf{RF/RNF} & \textbf{Use Case} & \textbf{User Story} & \textbf{Acceptance Criteria} & \textbf{Mockup} & \textbf{Note} \\
\midrule
\endhead

RF-01 & UC-01 & US-01 & AC-US01-01/02/03 & MU-Fig.3 & Catalogo edifici. \\
RF-02 & UC-13 & US-02 & AC-US02-01/02 & MU-Fig.2 & Vista cartografica. \\
RF-03 & UC-02 & US-03 & AC-US03-01 & MU-Fig.4 & Catalogo interventi. \\
RF-04 & UC-02 & US-03 & AC-US03-02/03 & MU-Fig.4 & Ricerca e filtro interventi. \\
RF-05 & UC-11 & US-04 & AC-US04-01/02/03 & MU-Fig.9 & Import CSV in due fasi. \\
RF-06 & UC-03 & US-05 & AC-US05-01/02/03 & MU-Fig.5 & Configurazione scenario. \\
RF-07 & UC-04 & US-06 & AC-US06-01/02/03 & MU-Fig.5 & Vincoli dello scenario. \\
RF-08 & UC-03 & US-05 & AC-US05-03 & MU-Fig.5 & Selezione candidati. \\
RF-09 & UC-05 & US-07 & AC-US07-01/02/03 & MU-Fig.6 & Ottimizzazione. \\
RF-10 & UC-06 & US-07 & AC-US07-01 & MU-Fig.6 & Visualizzazione risultati. \\
RF-11 & UC-06 & US-08 & AC-US08-01/02/03 & MU-Fig.6 & Spiegazioni deterministiche. \\
RF-12 & UC-12 & Da definire & Da definire & Da definire & Audit: consultazione storica, priorità media. \\
RF-13 & UC-07 & US-09 & AC-US09-01/02 & MU-Fig.7 & Salvataggio A/B. \\
RF-14 & UC-07 & US-10 & AC-US10-01/02 & MU-Fig.7 & Confronto scenari. \\
RF-15 & UC-08 & US-11 & AC-US11-01/02 & MU-Fig.8 & Sensibilità sul budget. \\
RF-16 & UC-09 & US-12 & AC-US12-01/02 & MU-Fig.5 & Proposta AI opzionale. \\
RF-17 & UC-09, UC-10 & US-13 & AC-US13-01/02/03 & MU-Fig.5 & Approvazione umana. \\
RF-18 & --- & US-14 & AC-US14-01/02 & MU-Fig.1 & Cambio tema. \\
\midrule
RNF-01 & --- & --- & --- & --- & Sezione 3, motivazione RNF-01. \\
RNF-02 & UC-05 & US-07 & AC-US07-02 & --- & Determinismo solver. \\
RNF-03 & UC-12 & Da definire & Da definire & --- & Audit. \\
RNF-04 & UC-06 & US-08 & AC-US08-01/02/03 & MU-Fig.6 & Interpretabilità. \\
RNF-05 & UC-05 & US-07 & AC-US07-03 & --- & Prestazioni. \\
RNF-06 & --- & --- & --- & --- & Compatibilità browser. \\
RNF-07 & --- & --- & --- & MU-Fig.1..9 & Accessibilità WCAG 2.1 AA. \\
RNF-08 & --- & --- & --- & --- & Manutenibilità / OpenAPI. \\
RNF-09 & UC-09, UC-10 & US-12, US-13 & AC-US12-01/02, AC-US13-01/02/03 & MU-Fig.5 & AI opzionale e confinata. \\
RNF-10 & UC-12 & Da definire & Da definire & --- & Snapshot immutabile. \\

\bottomrule
\end{longtable}

\textbf{Note:} gli elementi RF-12, RNF-03 e RNF-10 sono esplicitamente a priorità media e sono documentati in questa sede ma non ancora collegati a uno user flow dedicato. Saranno oggetto di aggiornamento dopo la baseline.
\section{Gestione delle modifiche dopo la baseline}

Questa sezione documenta le modifiche significative apportate alla baseline, riportandone una breve sintesi e il collegamento alle relative GitHub Issue e/o Pull Request. La cronologia completa resta tracciata nel repository GitHub.

\subsection*{Tipi di modifica}
\begin{itemize}
    \item \textbf{Req./UC/US/AC/MU Change:} modifica di un requisito, use case, user story, acceptance criteria o mockup.
    \item \textbf{Clarification:} rende esplicito un significato precedentemente ambiguo o incompleto.
\end{itemize}

\subsection*{Registro delle modifiche}

\begin{longtable}{p{1.6cm} p{1.7cm} p{4.5cm} p{2.2cm} p{3cm} p{2.5cm}}
\caption{Registro delle modifiche di TEPO successive alla baseline.}\\
\toprule
\textbf{ID} & \textbf{Tipo} & \textbf{Descrizione} & \textbf{Origine} & \textbf{Artefatti impattati} & \textbf{GitHub / Stato} \\
\midrule
\endfirsthead
\toprule
\textbf{ID} & \textbf{Tipo} & \textbf{Descrizione} & \textbf{Origine} & \textbf{Artefatti impattati} & \textbf{GitHub / Stato} \\
\midrule
\endhead

CH-TEPO-01 & Clarification & La soglia di prestazione dell'ottimizzazione è riferita a 30 opzioni candidate (non 26 come inizialmente indicato), in coerenza con l'estensione prevista del catalogo. & Team review & RNF\ref{rnf:prestazioni}; US\ref{us:ottimizza}; AC-US07-03 & \#12 --- DONE \\
CH-TEPO-02 & Req. Change & Aggiunta esplicita, fra i limiti dichiarati, dell'assenza di pianificazione multi-anno e di frontiera di Pareto. & Stakeholder (debriefing) & Sezione 1; RNF\ref{rnf:isolamento} & \#18 --- DONE \\
CH-TEPO-03 & Clarification & Chiarito che la componente AI può ricadere in modo automatico sul motore a regole qualora il provider esterno non sia configurato o non risponda. & Team review & RF\ref{rf:ai_proposta}; RNF\ref{rnf:ai_isolamento}; UC-09; AC-US12-02 & \#22 --- DONE \\
CH-TEPO-04 & MU Change & Ridisegno del pannello di proposta AI con evidenziazione degli interventi proposti come esclusi senza deselezione automatica. & User feedback interno & MU-Fig.5; US\ref{us:ai_appr}; AC-US13-01 & \#25 --- OPEN \\
CH-TEPO-05 & Req. Change & Introdotto requisito di snapshot immutabile delle revisioni dati (RNF\ref{rnf:snapshot}) per garantire la ricostruibilità dei run di audit. & Audit interno & RNF\ref{rnf:snapshot}; RF\ref{rf:storico}; UC-12 & \#28 --- OPEN \\
CH-TEPO-06 & Clarification & Chiarito che la pagina di sensibilità esegue uno sweep deterministico sul budget e non una simulazione Monte Carlo. & Team review & RF\ref{rf:sensibilita}; US\ref{us:sensibilita}; AC-US11-01 & \#31 --- DONE \\

\bottomrule
\end{longtable}

\textbf{Nota operativa.} Dopo la baseline, il D1 è un \emph{living artifact}: le modifiche significative vengono tracciate su GitHub e sottoposte a impact analysis. Gli artefatti successivi (D2, codice, test) sono indicati come impattati solo se già esistenti al momento della modifica. La cronologia completa è mantenuta nel repository, questa sezione ne riporta la sintesi.

\end{document}